\let\negmedspace\undefined
\let\negthickspace\undefined
\documentclass[journal,12pt,onecolumn]{IEEEtran}
\usepackage{cite}
\usepackage{amsmath,amssymb,amsfonts,amsthm}
\usepackage{amsmath}
\usepackage{algorithmic}
\usepackage{graphicx}
\usepackage{textcomp}
\usepackage{xcolor}
\usepackage{txfonts}
\usepackage{listings}
\usepackage{multicol}
\usepackage{enumitem}
\usepackage{mathtools}
\usepackage{gensymb}
\usepackage{comment}
\usepackage[breaklinks=true]{hyperref}
\usepackage{tkz-euclide} 
\usepackage{listings}
\usepackage{gvv}                                        
\usepackage[latin1]{inputenc}                                
\usepackage{color}                                            
\usepackage{array}                                            
\usepackage{longtable}                                       
\usepackage{calc}                                             
\usepackage{multirow}                                         
\usepackage{hhline}                                           
\usepackage{ifthen}                                           
\usepackage{lscape}
\usepackage{tabularx}
\usepackage{array}
\usepackage{float}
\usepackage{tikz}
\usepackage{multicol}
\usepackage{circuitikz}
\usepackage{amsmath}
\usetikzlibrary{patterns}
\newtheorem{theorem}{Theorem}[section]
\newtheorem{problem}{Problem}
\newtheorem{proposition}{Proposition}[section]
\newtheorem{lemma}{Lemma}[section]
\newtheorem{corollary}[theorem]{Corollary}
\newtheorem{example}{Example}[section]
\newtheorem{definition}[problem]{Definition}
\newcommand{\BEQA}{\begin{eqnarray}}
\newcommand{\EEQA}{\end{eqnarray}}
\newcommand{\define}{\stackrel{\triangle}{=}}
\theoremstyle{remark}
\newtheorem{rem}{Remark}

\begin{document}

\bibliographystyle{IEEEtran}
\vspace{3cm}

\title{2017-EE}
\author{EE24BTECH11020 -  Ellanti Rohith}
\maketitle

\renewcommand{\thefigure}{\theenumi}
\renewcommand{\thetable}{\theenumi}

\begin{enumerate}
\item Consider the differential equation 
\begin{align*}
\brak{t^2 - 81} \frac{dy}{dt} + 5t \, y = \sin\brak{t} \quad \text{with } y\brak{1} = 2\pi.
\end{align*}

There exists a unique solution for this differential equation when $ t$  belongs to the interval

\hfill{[GATE 2017]} \begin{enumerate}
    \begin{multicols}{2}
        \item $\brak{-2, 2}$
        \item $\brak{-10, 10}$
        \item $\brak{-10, 2}$
        \item $\brak{0, 10}$
    \end{multicols}
\end{enumerate}
\item Consider the line integral 
\begin{align*}
I = \int_C \brak{x^2 + iy} \, dz,
\end{align*}
where $ z = x + iy $. The line $ C $ is shown in the figure below.

\begin{center}
    
\resizebox{0.2\textwidth}{!}{%
\begin{circuitikz}
\tikzstyle{every node}=[font=\LARGE]
\draw  (6.25,10.5) circle (1.75cm);
\draw  (6.25,10.5) circle (0.5cm);
\draw  (6,11.75) rectangle (6.5,11.5);
\draw  (7.5,10.75) rectangle (7.75,10);
\draw  (6,9.5) rectangle (6.5,9.25);
\draw  (4.75,10.75) rectangle (5,10);
\end{circuitikz}
}%

\end{center}

The value of $ I $ is

\hfill{[GATE 2017]} \begin{enumerate}
    \begin{multicols}{2}
        \item $\dfrac{1}{2} i$\\
        \item  $\dfrac{2}{3} i$
        \item $\dfrac{3}{4} i$\\
        \item $\dfrac{4}{5} i$
    \end{multicols}
\end{enumerate}
\item Two passive two-port networks are connected in cascade as shown in the figure. A voltage source is connected at port 1.

\begin{center}
  \begin{figure}[!ht]
\centering
\resizebox{0.7\textwidth}{!}{%
\begin{circuitikz}
\tikzstyle{every node}=[font=\large]
\draw [ line width=1.3pt](6.25,11.75) to[sinusoidal voltage source, sources/symbol/rotate=auto] (6.25,8.75);
\draw [ line width=1.3pt](6.25,8.75) to[short] (7.75,8.75);
\draw [ line width=1.3pt](6.25,11.75) to[short] (7.75,11.75);
\draw [ line width=1.3pt](7.75,8.5) to[short] (7.75,12.5);
\draw [ line width=1.3pt](7.75,8.25) to[short] (7.75,9);
\draw [ line width=1.3pt ] (7.75,12.5) rectangle (13.75,8.25);
\draw [ line width=1.3pt ] (13.75,12) rectangle (17.5,8.75);
\draw [ line width=1.3pt ] (17.5,12.5) rectangle (24,8.5);
\draw [line width=1.3pt, short] (24,12) -- (25.5,12);
\draw [line width=1.3pt, short] (24,8.75) -- (25.5,8.75);
\draw [ line width=1.3pt](14,12.75) to[short] (15,12.75);
\draw [ line width=1.3pt](15,12.75) to[short] (14.5,13);
\draw [ line width=1.3pt](15,12.75) to[short] (14.5,12.5);
\node [font=\large] at (16,12.75) {$I_2$};
\draw [ line width=1.3pt](24,12.75) to[short] (25,12.75);
\draw [ line width=1.3pt](25,12.75) to[short] (24.5,12.5);
\draw [ line width=1.3pt](25,12.75) to[short] (24.5,13);
\node [font=\large] at(26,12.75)  {$I_3$};

\draw [ line width=1.3pt](5.25,12.5) to[short] (6.25,12.5);
\draw [ line width=1.3pt](6.25,12.5) to[short] (6,12.25);
\draw [ line width=1.3pt](6.25,12.5) to[short] (6,12.75);
\node [font=\large] at (7,12.5) {$I_1$};

\node [font=\large] at (10.75,10.25) {Two-Port Network 1};
\node [font=\large] at (21,10.25) {Two-Port Network 2};

\node [font=\large] at(25,10.5)  {$V_2$};
\node [font=\large] at (5.5,10.5) {$V_1$};
\node [font=\large] at (15.5,10.5) {$V_3$};

\node [font=\large] at (7,8) {Port 1};
\node [font=\large] at (15.5,8) {Port 2};
\node [font=\large] at (25,8) {Port 3};
\node [font=\large] at (5.5,11.2) {+};
\node [font=\large] at (15.5,11.5) {+};
\node [font=\large] at (24.75,11.5) {+};
\node [font=\large] at (24.75,9.5) {-};
\node [font=\large] at (15.5,9.5) {-};
\node [font=\large] at (5.5,9.7) {-};
\end{circuitikz}
}%


\end{figure}

\end{center}

Given
\begin{align*}
    V_1 &= A_1 V_2 + B_1 I_2\\ I_1 &= C_1 V_2 + D_1 I_2\\
    V_2 &= A_2 V_3 + B_2 I_3 \\ I_2 &= C_2 V_3 + D_2 I_3
\end{align*}
    



where  $A_1, B_1, C_1, D_1, A_2, B_2, C_2$ and $ D_2$ are the generalized circuit constants. If the Thevenin equivalent circuit at port 3 consists of a voltage source  $V_T$ and an impedance $ Z_T$, connected in series, then

\hfill{[GATE 2017]} \begin{enumerate}
    \begin{multicols}{2}
        \item  $V_T = \dfrac{V_1}{A_1 + A_2}, \quad Z_T = \dfrac{A_1 B_2 + B_1 D_2}{A_1 A_2}$\\
        \item  $V_T = \dfrac{V_1}{A_1 A_2 + B_1 C_2}, \quad Z_T = \dfrac{A_1 B_2 + B_1 D_2}{A_1 A_2 + B_1 C_2} $
        \item$ V_T = \dfrac{V_1}{A_1 + A_2}, \quad Z_T = \dfrac{A_1 A_2 + B_1 C_2}{A_1 A_2 + B_1 D_2}$\\
        \item $ V_T = \dfrac{V_1}{A_1 A_2 + B_1 C_2}, \quad Z_T = \dfrac{A_1 A_2 + B_1 C_2}{A_1 A_2 + B_1 C_2}$
    \end{multicols}
\end{enumerate}

\item Let a causal LTI system be characterized by the following differential equation, with initial rest condition
\begin{align*}
   \frac{d^2 y}{dt^2} + 7 \frac{dy}{dt} + 10 y\brak{t} = 4x\brak{t} + 5 \frac{dx\brak{t}}{dt} 
\end{align*}

where $x\brak{t}$ and  $y\brak{t}$ are the input and output respectively. The impulse response of the system is $ u\brak{t}$ is the unit step function

\hfill{[GATE 2017]} \begin{enumerate}
    \begin{multicols}{2}
        \item  $2e^{-2t} u\brak{t} - 7e^{-5t} u\brak{t} $
        \item  $-2e^{-2t} u\brak{t} + 7e^{-5t} u\brak{t} $
        \item $ 7e^{-2t} u\brak{t} - 2e^{-5t} u\brak{t} $
        \item $-7e^{-2t} u\brak{t} + 2e^{-5t} u\brak{t} $
    \end{multicols}
\end{enumerate}
\item Let the signal
\begin{align*}
  x\brak{t} = \sum_{k=-\infty}^{+\infty} \brak{-1}^k \delta \brak{ t - \frac{k}{2000} }
\end{align*}

be passed through an LTI system with frequency response $H\brak{\omega}$, as given in the figure below.
\begin{figure}[!ht]
\centering
\resizebox{0.27\textwidth}{!}{%
\begin{circuitikz}
\tikzstyle{every node}=[font=\normalsize]
\draw [->, >=Stealth] (1.25,9.25) -- (5,9.25);
\draw [->, >=Stealth] (1.25,9.25) -- (1.25,13.25);
\node [font=\normalsize, rotate around={90:(0,0)}] at (0.75,11.25) {Stress};
\node [font=\normalsize] at (3.25,9) {Strain};
\draw [->, >=Stealth] (1.25,9.25) -- (2.5,10.5);
\draw [->, >=Stealth] (4.5,12.5) -- (3.25,11.25);
\draw (2.5,10.5) to[short] (3.25,11.25);
\end{circuitikz}
}%

\end{figure}



The Fourier series representation of the output is given as\\

\hfill{[GATE 2017]} \begin{enumerate}
 
        \item $ 4000 + 4000 \cos\brak{2000 \pi t} + 4000 \cos\brak{4000 \pi t} $   \item$ 2000 + 2000 \cos\brak{2000 \pi t} + 2000 \cos\brak{4000 \pi t}$
        \item $ 2000 \cos\brak{2000 \pi t} $
       \item $ 4000 \cos\brak{2000 \pi t}$\\
    
\end{enumerate}

\item In the system whose signal flow graph is shown in the figure,$ U_1\brak{s}$ and  $U_2\brak{s}$ are inputs. The transfer function $\dfrac{Y\brak{s}}{U_1\brak{s}}$  is
\hfill{[GATE 2017]} 
\begin{center}
\resizebox{0.6\textwidth}{!}{%
\begin{circuitikz}
\tikzstyle{every node}=[font=\large]
\draw (2.5,13) to[short] (2.5,6);
\draw (2.5,6) to[short] (11.25,6);
\draw (11.25,13) to[short] (11.25,6);
\draw [color={rgb,255:red,0; green,128; blue,0}](2.5,6.5) to[short] (11.25,6.5);
\draw [ color={rgb,255:red,255; green,0; blue,0}, dashed] (3,12.75) .. controls (3,9.25) and (5.5,7) .. (11.25,10.5);
\draw [color={rgb,255:red,0; green,0; blue,255}, dashed] (3,12.75) .. controls (3.5,6) and (5.25,7.75) .. (11.5,7);
\draw [ color={rgb,255:red,128; green,0; blue,128}, dashdotted] (2.5,6) -- (11.25,9.75);
\node [font=\large, rotate around={90:(0,0)}] at (2.1,9.5) {$Cost\text{ } per\text{ } unit \rightarrow$};
\node [font=\large] at (7.75,5.5) {$Order\text{ } quantity \rightarrow$};
\node [font=\normalsize, color={rgb,255:red,255; green,0; blue,0}] at (9.25,10.75) {$Curve \, P1$};
\node [font=\normalsize, color={rgb,255:red,128; green,0; blue,128}] at (10.5,8.5) {$Curve \, P2$};
\node [font=\normalsize, color={rgb,255:red,0; green,0; blue,255}] at (8.75,8) {$Curve \, P3$};
\node [font=\normalsize, color={rgb,255:red,0; green,128; blue,0}] at (5.75,7) {$Curve \, P4$};
\draw [->, >=Stealth] (6.5,7) -- (7.75,6.5);
\draw [->, >=Stealth] (9,7.75) -- (9.5,7.25);
\draw [->, >=Stealth] (9.25,10.5) -- (9.5,9.75);
\draw [->, >=Stealth] (10.5,8.75) -- (10,9.25);
\end{circuitikz}
}%
\end{center}
\begin{enumerate}
    \begin{multicols}{2}
        \item $\dfrac{k_1}{JL s^2 + JR s + k_1 k_2}$
        \item  $\dfrac{k_1}{JL s^2 - JR s - k_1 k_2}$
        \item  $\dfrac{k_1 - U_2 \brak{R + sL}}{JL s^2 + \brak{JR - U_2 L} s + k_1 k_2 - U_2 R}$
        \item  $\dfrac{k_1 - U_2 \brak{sL - R}}{JL s^2 - \brak{JR + U_2 L} s - k_1 k_2 + U_2 R}$
    \end{multicols}
\end{enumerate}
\item The transfer function of the system is given by:

\begin{align*}
x\brak{t} &= \myvec{ 1 & 2 \\ 2 & 0} + \myvec {1 \\ 2 } u\brak{t}  \\
y\brak{t} &= \myvec{ 1 & 0} x\brak{t}
\end{align*}
The options for the transfer function are:\hfill{[GATE 2017]} 
\begin{multicols}{2}
\begin{enumerate}
    \item $\dfrac{\brak{s + 2}}{s^2 - 2s - 2}$\\
    \item  $\dfrac{\brak{s - 2}}{s^2 + s - 4}$
    \item  $\dfrac{\brak{s + 4}}{s^2 + s - 4}$\\
    \item  $\dfrac{\brak{s + 4}}{s^2 - s - 4}$
\end{enumerate}     
\end{multicols}

\item The load shown in the figure is supplied by a 400 $V$ (line-to-line), 3-phase source RYB sequence. The load is balanced and inductive, drawing 3464 $VA$. When the switch $S$ is in position $ N $, the three wattmeters $W_1 $,  $W_2$, and $ W_3 $ read 577.35 $W$ each. If the switch is moved to position $Y$, the readings of the wattmeters in watts will be:
\hfill{[GATE 2017]} 

\begin{center}
\resizebox{0.7\textwidth}{!}{%
\begin{circuitikz}
\tikzstyle{every node}=[font=\LARGE]
\draw [->, >=Stealth] (2.5,10.5) -- (3.75,10.5);
\draw [->, >=Stealth] (2.5,10.125) -- (3.75,10.125);
\draw [->, >=Stealth] (2.5,9.75) -- (3.75,9.75);
\draw [->, >=Stealth] (10,10.5) -- (11.25,10.5);
\draw [->, >=Stealth] (10,10.125) -- (11.25,10.125);
\draw [->, >=Stealth] (10,9.75) -- (11.25,9.75);
\draw [short] (6.25,8.75) .. controls (2.75,11.75) and (4.75,12.5) .. (6.25,8.75);
\draw [short] (13.75,8.75) .. controls (10,14) and (13,12.25) .. (13.64,8.9);
\node [font=\LARGE] at (9.5,10.25) {U};
\node [font=\normalsize] at (1.5,8.25) {Trailing Edge With Finite Angle};
\node [font=\normalsize] at (9.5,8.25) {Trailing Edge With cusp};
\node [font=\LARGE] at (1.75,10.25) {U};
\end{circuitikz}
}%
\end{center}
\item The approximate transfer characteristic for the circuit shown below with an ideal operational amplifier and diode is as follows:

\begin{center}
\resizebox{0.3\textwidth}{!}{%
\begin{circuitikz}
\tikzstyle{every node}=[font=\large]
\fill[lightgray] (6.25,8.75) -- (7.75,10.25) -- (6.25,11.75) --(4.5,10) --cycle;
\draw [black, line width=1.2pt](6.25,11.75) to[short] (6.25,8.75);
\draw [black, line width=1.2pt](6.25,11.75) to[short] (7.75,10.25);
\draw [black, line width=1.2pt](6.25,8.75) to[short] (7.75,10.25);
\draw [black, line width=1.2pt](6.25,11.75) to[short] (4.5,10);
\draw [black, line width=1.2pt](4.5,10) to[short] (6.25,8.75);
\draw [black, line width=1.2pt, dashed] (4.5,10) -- (6.55,10.3);
\draw [black, line width=1.2pt, dashed] (6.25,11.75) -- (6.55,10.3);
\draw [black, line width=1.2pt, dashed] (6.55,10.3) -- (7.75,10.25);



\end{circuitikz}
}%
\end{center}\hfill{[GATE 2017]}
\begin{multicols}{2}
     \begin{enumerate}
    \item 
\resizebox{0.25\textwidth}{!}{%
\begin{circuitikz}
\tikzstyle{every node}=[font=\LARGE]
\draw [->, >=Stealth] (2.5,8.25) -- (2.5,11.75);
\draw [->, >=Stealth] (2.5,8.25) -- (6.25,8.25);
\node [font=\normalsize] at (2.5,12) {$C_{L}$};
\node [font=\normalsize] at (6.25,8.75) {$C_{D}$};
\node [font=\normalsize] at (2.25,8) {$O$};
\draw [short] (2.5,10.75) .. controls (4,7.75) and (4.5,7.25) .. (5.75,10.75);
\end{circuitikz}
}%

 
     \item 
\resizebox{0.25\textwidth}{!}{%
\begin{circuitikz}
\tikzstyle{every node}=[font=\LARGE]
\draw [->, >=Stealth] (2.5,8.25) -- (2.5,11.75);
\draw [->, >=Stealth] (2.5,8.25) -- (6.25,8.25);
\node [font=\normalsize] at (2.5,12) {$C_{L}$};
\node [font=\normalsize] at (6.25,8.75) {$C_{D}$};
\node [font=\normalsize] at (2.25,8) {$O$};
\draw [short] (2.5,10.75) .. controls (5.25,9.75) and (5.75,9.5) .. (2.5,8.25);
\end{circuitikz}
}%
  \item 
\resizebox{0.25\textwidth}{!}{%
\begin{circuitikz}
\tikzstyle{every node}=[font=\LARGE]
\draw [->, >=Stealth] (2.5,8.25) -- (2.5,11.75);
\draw [->, >=Stealth] (2.5,8.25) -- (6.25,8.25);
\node [font=\normalsize] at (2.5,12) {$C_{L}$};
\node [font=\normalsize] at (6.25,8.75) {$C_{D}$};
\node [font=\normalsize] at (2.25,8) {$O$};
\draw [short] (2.5,8.25) .. controls (3.5,11.75) and (4.5,11.25) .. (5.25,8.25);
\end{circuitikz}
}%
  \item 
\resizebox{0.15\textwidth}{!}{%
\begin{circuitikz}
\tikzstyle{every node}=[font=\normalsize]

\filldraw[fill=gray!30, draw=black] 
    (5,8.25) -- (6.5,9) -- (8,7.75) -- (6.5,5.75) -- cycle;

\node at (6.5,8) {\textbf{R}};

\end{circuitikz}
}%
\end{enumerate}
\end{multicols}

   

    

\item The output expression for the Karnaugh map shown below is
\begin{align*}
  \begin{array}{|c|c|c|c|c|}\hline
    AB \backslash CD & 00 & 01 & 11 & 10 \\
    \hline
    00 & 0 & 0 & 0 & 0 \\\hline
    01 & 1 & 0 & 1 & 1 \\ \hline
    11 & 1 & 0 & 1 & 1 \\ \hline
    10 & 0 & 0 & 0 & 0 \\ \hline
\end{array}  
\end{align*}


\hfill{[GATE 2017]}
\begin{multicols}{4}
\begin{enumerate}
    \item $ \overline{B} \overline{D} + BCD $
    \item $ \overline{B} D + AB $
    \item $ \overline{B} D + ABC $
    \item $ B \overline{D} + ABC $
\end{enumerate}
\end{multicols}
\item A load is supplied by a 230 $V$, 50 $Hz$ source. The active power $ P $ and the reactive power $ Q $ consumed by the load are such that $1 \text{ kW} \leq P \leq 2 \text{ kW}$ and $ 1 \text{ kVAR} \leq Q \leq 2 \text{ kVAR}$. A capacitor connected across the load for power factor correction generates 1 $kVAR$ reactive power. The worst case power factor after power factor correction is
\hfill{[GATE 2017]} 
\begin{multicols}{4}
 \begin{enumerate}
    \item 0.447 lag
    \item 0.707 lag
    \item 0.894 lag
    \item 1
\end{enumerate}
\end{multicols}

\item The logical gate implemented using the circuit shown below where $ V_1$ and $ V_2 $ are inputs with (0 $V$ as digital 0 and 5 $V$ as digital 1) and $ V_{\text{OUT}} $ is the output, is \\ \newpage
\begin{figure}[!ht]
\centering
\resizebox{0.3\textwidth}{!}{%
\begin{circuitikz}
\tikzstyle{every node}=[font=\large]
\draw [ line width=1.2pt](3.75,11.75) to[short] (3.75,6.75);
\draw [ line width=1.2pt](3.75,6.75) to[short] (8.75,6.75);
\draw [ line width=1.2pt](8.75,11.75) to[short] (8.75,6.75);
\draw [ line width=1.2pt](3.75,11.75) to[short] (8.75,11.75);
\draw [line width=1.2pt, dashed] (5,11.75) -- (5,6.75);
\draw [line width=1.2pt, dashed] (7.5,11.75) -- (7.5,6.75);
\draw [line width=1.2pt, dashed] (3.75,10.5) -- (8.75,10.5);
\draw [line width=1.2pt, dashed] (3.75,8) -- (8.75,8);
\draw [ line width=1.2pt](6.25,11.75) to[short] (6.25,6.75);
\draw [ line width=1.2pt](3.75,9.25) to[short] (8.75,9.25);
\node [font=\large] at (4.35,7.25) {$\ast$};
\node [font=\large] at (7,9.9) {$\ast$};
\node [font=\large] at (8.1,8.5) {$\ast$};
\node [font=\large] at (8.25,11) {$\triangle$};
\node [font=\large] at (7,7.35) {$\triangle$};
\node [font=\large] at (5.7,9.8) {$\circ$};
\node [font=\large] at (8.1,7.25) {$\circ$};
\node [font=\large] at (7,8.5) {$\square$};
\node [font=\large] at (5.7,11) {$?$};
\end{circuitikz}
}%

\end{figure}
\hfill{[GATE 2017]}
\item The input voltage $V_{DC}$ of the buck-boost converter shown below varies from 32 $V$ to 72 $V$. Assume that all components are ideal, inductor current is continuous, and output voltage is ripple-free. The range of duty ratio $D$ of the converter for which the magnitude of the steady-state output voltage remains constant at 48 $V$ is
\hfill{[GATE 2017]}
 \begin{multicols}{2}
    \begin{enumerate}
 \item \begin{tikzpicture}
            \fill[gray] (0,0) -- (-0.5,1) -- (1,1) -- (1.5,0)-- cycle;
              \fill[gray] (0,0) -- (-0.5,-1) -- (1,-1) -- (1.5,0) -- cycle;
        \end{tikzpicture}        
        \item \begin{tikzpicture}
            \fill[gray] (0,0) -- (1,1.73) -- (2,0) -- cycle;
        \end{tikzpicture}
       \item \begin{tikzpicture}
            \fill[gray] (0,0) -- (2.2,0) -- (1.85,1.5) -- (0.35,1.5) -- cycle; 
            
        \end{tikzpicture}
        \item  \begin{tikzpicture}
            \fill[gray] (0,0) rectangle (2,2);
        \end{tikzpicture}
        
        
    \end{enumerate}
    \end{multicols}

\begin{multicols}{4}
\begin{enumerate}
    \item $ \dfrac{2}{5} \leq D \leq \dfrac{3}{5} $
    \item $ \dfrac{2}{3} \leq D \leq \dfrac{3}{4} $
    \item $ 0 \leq D \leq 1 $
    \item $ \dfrac{1}{3} \leq D \leq \dfrac{2}{3} $
\end{enumerate}

\end{multicols}




\end{enumerate}


\end{document}
